\subsection{Versiones de salida múltiple: diez pasos}
\label{subsec:lineales-multi}

En \texttt{OLS\_10.py} y \texttt{ridge\_10.py} se construyeron ventanas de 22 registros, con una historia nominal de 12 y un horizonte de diez. Cada respuesta de entrenamiento es un vector de diez valores consecutivos de generación posteriores a la fila de entrada. El vector pronosticado se obtiene usando los predictores de la última fila de la historia. OLS utiliza los cinco predictores de calendario y estandarización; Ridge añade el rezago y conserva la selección de $\alpha$ con tres particiones, sin estandarizar.

La construcción de las respuestas recorre las 12 filas de historia, pero extrae sus valores de la ventana completa de 22 registros. En particular, la respuesta asociada a la última fila histórica contiene exactamente los diez valores que después se usan para evaluar esa ventana. Además, esa misma fila de entrada participa en el ajuste y en la predicción. Existe, por tanto, fuga de información de las respuestas de evaluación hacia el entrenamiento. Las salidas conservadas describen este procedimiento exploratorio y no una prueba independiente de pronóstico a diez pasos.

Los programas calculan MAE y $R^2$ promediados entre las diez salidas. El valor denominado RMSE es la raíz del error cuadrático medio agregado entre observaciones y salidas; no es el promedio de diez RMSE calculados por separado. Las gráficas \texttt{prediccion\_ols\_multioutput.png} y \texttt{prediccion\_ridge\_multioutput.png} muestran únicamente el primer paso.
